\section{What Computers Do ?}

At a high level, a computer is a programmable electronic system that receives data, processes them according to a set of instructions, stores information, and produces results. Although modern computers perform an enormous variety of tasks, from running websites and analysing scientific data to controlling industrial systems, these activities can ultimately be understood as different forms of information processing. \\


The fundamental operation of a computer is the execution of instructions. Software provides these instructions in the form of programs, while hardware provides the physical components needed to execute them. A program might instruct a computer to perform calculations, compare values, move data between locations, communicate over a network, or respond to user input.\\

The central processing unit (CPU) is the centre of most computers. The CPU repeatedly retrieves instructions from memory, interprets them, and performs the required operations. These operations are ultimately carried out using electronic circuits that manipulate binary values—represented conceptually as zeros and ones. Simple operations such as addition, comparison, and logical operations can be combined into increasingly complex sequences, allowing software to perform sophisticated tasks. \\

Computers also rely heavily on memory and storage. Main memory, or RAM, provides relatively fast temporary storage for the programs and data currently being used. Persistent storage, such as solid-state drives, retains data when the computer is switched off. The operating system manages these resources and provides common services through which applications can interact with hardware without having to directly control each physical component. \\

Input and output allow a computer to interact with the outside world. Input can come from devices such as keyboards, cameras, sensors, or network connections, while output can be presented through displays, speakers, actuators, or other connected systems. Internally, hardware components communicate through electronic interconnections and communication protocols, allowing data to move between processor, memory, storage, and peripheral devices. \\

At a higher level, the operating system coordinates these activities. Manage processes, memory, files, devices, and access to system resources. Applications such as web browsers, word processors and games operate above this layer, using operating-system services and software libraries rather than directly controlling the underlying hardware in most cases.\\

Modern computers also perform substantial work outside the individual machine. Through computer networks and the Internet, a device can exchange data with other computers and access remote services. In cloud computing, for example, an application may rely on processing and storage provided by large remote data centres rather than performing every operation locally.\\

The key idea is therefore that computers transform information through a hierarchy of abstraction. At the lowest level, electronic circuits manipulate binary states. These circuits implement processor instructions, which are combined into machine-level programs. Operating systems and software libraries provide higher-level abstractions, and applications use those abstractions to implement useful functionality.\\

Consequently, what appears to a user as a complex operation—such as editing a photograph, searching the Internet, or playing a video—is ultimately a very large sequence of relatively simple operations on data. The power of computing comes not primarily from any single operation, but from the speed, precision, and scale with which computers can execute these operations and combine them into increasingly sophisticated systems.